\chapter{Phase 4 --- Manuscript and Thesis}
\label{chap:phase4}

\begin{milestone}
\textbf{Status:} planned.\\
\textbf{Targets:} (i) \emph{Planetary Science Journal} (PSJ) submission
by 2026-06; (ii) dissertation chapter draft at the Institute of
Science Tokyo by 2026-08, with public defence window 2026-09.\\
\textbf{Deliverable:} \code{notebooks/phase4\_manuscript/} holding one
notebook per manuscript figure, each emitting a named PDF to
\code{output/figures/manuscript/}.
\end{milestone}

\section{Manuscript outline}

\begin{table}[H]
\centering
\small
\begin{tabularx}{\linewidth}{l l Y}
\toprule
Section & Length & Content \\
\midrule
Abstract        & 250 words  & Four novelty claims from \cref{chap:novelty}. \\
Introduction    & 1.5 p      & Volatiles context; gap in the literature. \\
Numerical model & 2 p        & \crefrange{eq:heat1d}{eq:bottom_bc}; grid; spin-up. \\
Phase-1 validation & 3 p      & Apollo, CE-4, ChaSTE (\cref{tab:phase1-rmse}). \\
Topographic coupling & 2 p    & \crefrange{eq:cos_incidence}{eq:Qtotal}; slope reruns. \\
Ice-coupled closure  & 2 p    & \crefrange{eq:k_mix}{eq:Tstab}; fixed-point loop. \\
Polar tile          & 3 p     & \cref{fig:phase3-loop} result; stability-depth map. \\
Discussion          & 2 p     & Sensitivity to $\phi_0$; comparison with \citet{paige2010}. \\
Conclusions         & 1 p     & Four-point novelty summary. \\
Supplementary       & --      & Every notebook and dataset via Zenodo DOI. \\
\bottomrule
\end{tabularx}
\caption{Manuscript outline targeting \emph{Planetary Science Journal}.}
\label{tab:psj-outline}
\end{table}

\section{Figure bundle}

\begin{table}[H]
\centering
\small
\begin{tabularx}{\linewidth}{l Y l}
\toprule
Figure & Content & Source notebook \\
\midrule
1 & Pipeline flow (\cref{fig:pipeline-flow}) & TikZ, this manual \\
2 & Apollo 15/17 mean-$T$ profile (\cref{fig:a15-meanT,fig:a17-meanT}) & \code{01\_apollo\_validation} \\
3 & A15 diurnal cycle (\cref{fig:a15-diurnal}) & \code{01\_apollo\_validation} \\
4 & CE-4 + ChaSTE flat-surface (\cref{fig:change4-val,fig:chaste-val}) & \code{01\_apollo\_validation} \\
5 & DEM + horizon + shadow panel (Phase~2) & \code{01\_slope\_horizon\_shadow} \\
6 & ChaSTE slope-corrected re-run (Phase~2) & \code{02\_chaste\_slope\_rerun} \\
7 & Polar tile $T_{\mathrm{bol}}$ vs.\ Diviner (Phase~3) & \code{01\_polar\_tile\_validation} \\
8 & Ice-stability depth map (Phase~3) & \code{02\_ice\_stability\_depth} \\
9 & $\phi_0$ sensitivity sweep (Phase~3) & \code{02\_ice\_stability\_depth} \\
\bottomrule
\end{tabularx}
\caption{Manuscript figure bundle.  Every figure is regenerated by
  \emph{Run All} on its source notebook.}
\label{tab:psj-figures}
\end{table}

\section{Style conventions}

\begin{itemize}
  \item \textbf{Colour palette.}  Discrete model $\to$
        \code{\#1f77b4} (blue), Hayne~(2017) $\to$ \code{\#d62728}
        (red), observations $\to$ black.  Per-site styling uses the
        \code{matplotlib} \code{tab10} family in the order Apollo~15,
        Apollo~17, Chang'E-4, ChaSTE.
  \item \textbf{Fonts.}  All manuscript figures use DejaVu Sans at
        \SI{9}{pt} body / \SI{11}{pt} axis labels, matched by the
        \code{lunar.plotting.set\_psj\_style()} helper.
  \item \textbf{File format.}  PDF primary (vector, $300\,\mathrm{DPI}$
        raster fallback embedded for preview); PNG secondary for the
        Overleaf previewer.
\end{itemize}

\section{Writing schedule}

\begin{table}[H]
\centering
\small
\begin{tabularx}{\linewidth}{l l Y}
\toprule
Milestone & Target date & Gate \\
\midrule
Phase 2 complete             & 2026-04-30 & ChaSTE bias $\le\SI{10}{\kelvin}$. \\
Phase 3 polar tile complete  & 2026-05-31 & Diviner RMS $\le\SI{5}{\kelvin}$. \\
Manuscript \S2--3 draft      & 2026-06-15 & Numerical model + Phase-1. \\
Manuscript \S4--5 draft      & 2026-06-30 & Topography + ice closure. \\
Manuscript \S6--7 draft      & 2026-07-15 & Polar tile + discussion. \\
Internal review              & 2026-07-31 & Supervisor + two co-authors. \\
PSJ submission               & 2026-08-15 & arXiv posting same day. \\
Dissertation chapter draft   & 2026-08-31 & Combined Phase-1..4 volume. \\
Final thesis submission      & 2026-09-15 & ISCT defence window. \\
\bottomrule
\end{tabularx}
\caption{Writing and submission schedule.  Every date is absolute; the
  Phase-2/3 gates are hard prerequisites for the manuscript drafts.}
\label{tab:schedule}
\end{table}

\section{Reproducibility artefact}

At submission time, a Zenodo-deposited artefact will accompany the
manuscript, containing:
\begin{enumerate}
  \item a frozen git tag (\code{v1.0-psj-submission}) of the
        \repo{} repository;
  \item a self-contained Docker image that runs every notebook
        end-to-end;
  \item the SHA-256 manifest of every input file (DEM, Diviner L2
        granules, SPICE kernels, HFE tables) together with the
        public-archive URLs.
\end{enumerate}
This satisfies the PSJ ``Data Availability'' policy in full and
supports the reproducibility claim of \cref{sec:novelty-reproducibility}.

\section{Closing remarks}

The four phases documented in this manual form a single coherent
argument: a physically transparent 1-D kernel, stress-tested against
every in-situ record, can be extended to a pixel-level polar map
with an ice-coupled conductivity closure, and the combination is
sufficient to make quantitative predictions about the depth of
stable water ice at the lunar south pole.  The remaining work is
chiefly engineering, not physics.  The finish line is \num{18}~months
of focused execution away.
